\begin{table}[htbp]
\centering
\caption{Standardized Effect Sizes}
\label{tab:sde}
\small
\begin{tabular}{lcccccc}
\hline\hline
\multicolumn{7}{l}{\textbf{Panel A: Pooled}} \\
\hline
Outcome & $\hat{\beta}$ & SE & SD($Y$) & SDE & SE(SDE) & Classification \\
\hline
Log catches (all) & -0.306 & 0.196 & 2.49 & -0.123 & 0.079 & Moderate negative \\
DDD (EU $\times$ demersal $\times$ post) & -1.226 & 0.411 & 2.13 & -0.576 & 0.193 & Large negative \\
\hline
\multicolumn{7}{l}{\textbf{Panel B: Heterogeneous (by species group)}} \\
\hline
Log catches (demersal) & -1.239 & 0.453 & 1.79 & -0.691 & 0.252 & Large negative \\
Log catches (pelagic) & 0.034 & 0.165 & 2.22 & 0.015 & 0.075 & Small positive \\
\hline\hline
\end{tabular}
\begin{tablenotes}[flushleft]
\footnotesize
\item \textit{Notes:} \textbf{Country:} European Union (15 member states with significant Atlantic fisheries: BE, DE, DK, EE, ES, FI, FR, IE, LT, LV, NL, PL, PT, SE, UK). \textbf{Research question:} Does the EU Landing Obligation (discard ban) reduce total catches through choke-species constraints in mixed fisheries, and does this effect differ between single-species pelagic and mixed demersal fisheries? \textbf{Policy mechanism:} The Landing Obligation (Regulation 1380/2013, Article 15) requires all catches of regulated species to be landed and counted against quotas, replacing the prior practice of discarding unwanted fish at sea; in mixed demersal fisheries, low-quota bycatch species become binding constraints that force vessels to stop fishing before exhausting target-species quotas (the ``choke species'' problem). \textbf{Outcome definition:} Log total catches (tonnes) by country $\times$ species group $\times$ year, from Eurostat fish\_ca\_main. \textbf{Treatment:} Binary; Landing Obligation activation staggered by species group (pelagic 2015, demersal 2016, other 2019). \textbf{Data:} Eurostat fisheries statistics (fish\_ca\_main), 2000--2024, country $\times$ species group $\times$ year; 1,252 observations across 51 units (45 treated EU, 6 non-EU control). \textbf{Method:} TWFE with unit and year fixed effects; triple-difference (EU $\times$ demersal $\times$ post) for mechanism test; standard errors clustered by country. \textbf{Sample:} EU member states with Atlantic fisheries plus Norway and Iceland as controls; species classified into Landing Obligation treatment cohorts based on phased implementation schedule. SDE $= \hat{\beta} / \text{SD}(Y)$ where SD($Y$) is the pre-treatment standard deviation. Classification refers to magnitude, not statistical significance: Large ($|$SDE$| > 0.15$), Moderate ($0.05$--$0.15$), Small ($0.005$--$0.05$), Null ($< 0.005$). 
\end{tablenotes}
\end{table}
